% Options for packages loaded elsewhere
% Options for packages loaded elsewhere
\PassOptionsToPackage{unicode}{hyperref}
\PassOptionsToPackage{hyphens}{url}
%
\documentclass[
  ignorenonframetext,
  aspectratio=169,
  russian,
  aspectratio=169]{beamer}
\newif\ifbibliography
\usepackage{pgfpages}
\setbeamertemplate{caption}[numbered]
\setbeamertemplate{caption label separator}{: }
\setbeamercolor{caption name}{fg=normal text.fg}
\beamertemplatenavigationsymbolshorizontal
% remove section numbering
\setbeamertemplate{part page}{
  \centering
  \begin{beamercolorbox}[sep=16pt,center]{part title}
    \usebeamerfont{part title}\insertpart\par
  \end{beamercolorbox}
}
\setbeamertemplate{section page}{
  \centering
  \begin{beamercolorbox}[sep=12pt,center]{section title}
    \usebeamerfont{section title}\insertsection\par
  \end{beamercolorbox}
}
\setbeamertemplate{subsection page}{
  \centering
  \begin{beamercolorbox}[sep=8pt,center]{subsection title}
    \usebeamerfont{subsection title}\insertsubsection\par
  \end{beamercolorbox}
}
% Prevent slide breaks in the middle of a paragraph
\widowpenalties 1 10000
\raggedbottom
\AtBeginPart{
  \frame{\partpage}
}
\AtBeginSection{
  \ifbibliography
  \else
    \frame{\sectionpage}
  \fi
}
\AtBeginSubsection{
  \frame{\subsectionpage}
}
\usepackage{iftex}
\ifPDFTeX
  \usepackage[T1]{fontenc}
  \usepackage[utf8]{inputenc}
  \usepackage{textcomp} % provide euro and other symbols
\else % if luatex or xetex
  \usepackage{unicode-math} % this also loads fontspec
  \defaultfontfeatures{Scale=MatchLowercase}
  \defaultfontfeatures[\rmfamily]{Ligatures=TeX,Scale=1}
\fi
\usepackage{lmodern}

\usefonttheme{serif} % use mainfont rather than sansfont for slide text
\ifPDFTeX\else
  % xetex/luatex font selection
  \setmainfont[]{DejaVu Serif}
  \setsansfont[]{DejaVu Sans}
  \setmonofont[]{DejaVu Sans Mono}
\fi
% Use upquote if available, for straight quotes in verbatim environments
\IfFileExists{upquote.sty}{\usepackage{upquote}}{}
\IfFileExists{microtype.sty}{% use microtype if available
  \usepackage[]{microtype}
  \UseMicrotypeSet[protrusion]{basicmath} % disable protrusion for tt fonts
}{}


\usepackage{longtable,booktabs,array}
\usepackage{calc} % for calculating minipage widths
\usepackage{caption}
% Make caption package work with longtable
\makeatletter
\def\fnum@table{\tablename~\thetable}
\makeatother
\usepackage{graphicx}
\makeatletter
\newsavebox\pandoc@box
\newcommand*\pandocbounded[1]{% scales image to fit in text height/width
  \sbox\pandoc@box{#1}%
  \Gscale@div\@tempa{\textheight}{\dimexpr\ht\pandoc@box+\dp\pandoc@box\relax}%
  \Gscale@div\@tempb{\linewidth}{\wd\pandoc@box}%
  \ifdim\@tempb\p@<\@tempa\p@\let\@tempa\@tempb\fi% select the smaller of both
  \ifdim\@tempa\p@<\p@\scalebox{\@tempa}{\usebox\pandoc@box}%
  \else\usebox{\pandoc@box}%
  \fi%
}
% Set default figure placement to htbp
\def\fps@figure{htbp}
\makeatother



\ifLuaTeX
\usepackage[bidi=basic,provide=*]{babel}
\else
\usepackage[bidi=default,provide=*]{babel}
\fi
\ifPDFTeX
\else
\babelfont{rm}[]{DejaVu Serif}
\fi
% get rid of language-specific shorthands (see #6817):
\let\LanguageShortHands\languageshorthands
\def\languageshorthands#1{}


\setlength{\emergencystretch}{3em} % prevent overfull lines

\providecommand{\tightlist}{%
  \setlength{\itemsep}{0pt}\setlength{\parskip}{0pt}}



 

\usepackage[]{csquotes}

\IfFileExists{plex-otf.sty}{
  %% Full TeXlive
  % \usepackage[%
  %   % math,
  %   RM={Scale=0.94},SS={Scale=0.94},SScon={Scale=0.94},TT={Scale=MatchLowercase,FakeStretch=0.9},DefaultFeatures={Ligatures=Common}
  % ]{plex-otf}
}{
  %% TinyTeX
  \usepackage{libertine}
}

%%% Load theme
% https://deic.uab.cat/~iblanes/beamer_gallery/
\IfFileExists{beamerthemegotham.sty}{
  %% Full TeXlive
  \usetheme{gotham}
  \gothamset{
    numbering=totalpagenumber,
    parttocframe default=off,
    sectiontocframe default=off,
    subsectiontocframe default=off,
  }
}{
  %% TinyTeX
  \usetheme{Madrid}
}
\makeatletter
\@ifpackageloaded{caption}{}{\usepackage{caption}}
\AtBeginDocument{%
\ifdefined\contentsname
  \renewcommand*\contentsname{Содержание}
\else
  \newcommand\contentsname{Содержание}
\fi
\ifdefined\listfigurename
  \renewcommand*\listfigurename{Список иллюстраций}
\else
  \newcommand\listfigurename{Список иллюстраций}
\fi
\ifdefined\listtablename
  \renewcommand*\listtablename{Список таблиц}
\else
  \newcommand\listtablename{Список таблиц}
\fi
\ifdefined\figurename
  \renewcommand*\figurename{Рисунок}
\else
  \newcommand\figurename{Рисунок}
\fi
\ifdefined\tablename
  \renewcommand*\tablename{Таблица}
\else
  \newcommand\tablename{Таблица}
\fi
}
\@ifpackageloaded{float}{}{\usepackage{float}}
\floatstyle{ruled}
\@ifundefined{c@chapter}{\newfloat{codelisting}{h}{lop}}{\newfloat{codelisting}{h}{lop}[chapter]}
\floatname{codelisting}{Список}
\newcommand*\listoflistings{\listof{codelisting}{Листинги}}
\makeatother
\makeatletter
\makeatother
\makeatletter
\@ifpackageloaded{caption}{}{\usepackage{caption}}
\@ifpackageloaded{subcaption}{}{\usepackage{subcaption}}
\makeatother

\usepackage{bookmark}
\IfFileExists{xurl.sty}{\usepackage{xurl}}{} % add URL line breaks if available
\urlstyle{same}
\hypersetup{
  pdftitle={Выполнение лабораторной работы №3},
  pdfauthor={Карпухин Клим},
  pdflang={ru-RU},
  hidelinks,
  pdfcreator={LaTeX via pandoc}}


\title{Выполнение лабораторной работы №3}
\subtitle{Оформление отчётов с помощью легковесного языка разметки
Markdown}
\author{Карпухин Клим}
\date{2026-03-01}

\begin{document}
\frame{\titlepage}


\section{Содержание}\label{ux441ux43eux434ux435ux440ux436ux430ux43dux438ux435}

\begin{frame}{Содержание}
\begin{enumerate}[<+->]
\tightlist
\item
  Информация о докладчике
\item
  Вводная часть и актуальность
\item
  Объект и предмет исследования
\item
  Научная новизна и практическая значимость
\item
  Цель и задачи
\item
  Материалы, методы и инструменты
\item
  Ход работы (этапы, скриншоты)
\item
  Результаты
\item
  Анализ и практическая значимость
\item
  Выводы
\end{enumerate}
\end{frame}

\section{Информация}\label{ux438ux43dux444ux43eux440ux43cux430ux446ux438ux44f}

\begin{frame}{Докладчик}
\phantomsection\label{ux434ux43eux43aux43bux430ux434ux447ux438ux43a}
\begin{columns}[c]
\begin{column}{0.65\linewidth}
\begin{itemize}[<+->]
\tightlist
\item
  \textbf{Карпухин Клим}
\item
  Российский университет дружбы народов
\item
  Email:
  \href{mailto:1032255580@rudn.ru}{\nolinkurl{1032255580@rudn.ru}}
\item
  Роли: студент (лабораторная работа по ОС/виртуализации)
\end{itemize}
\end{column}

\begin{column}{0.35\linewidth}

\includegraphics[width=0.9\linewidth,height=\textheight,keepaspectratio]{./image/dokladchik.png}
\end{column}
\end{columns}
\end{frame}

\section{Вводная
часть}\label{ux432ux432ux43eux434ux43dux430ux44f-ux447ux430ux441ux442ux44c}

\begin{frame}{Актуальность}
\phantomsection\label{ux430ux43aux442ux443ux430ux43bux44cux43dux43eux441ux442ux44c}
\begin{itemize}[<+->]
\tightlist
\item
  Markdown --- один из самых популярных языков разметки для оформления
  технической документации и отчётов.
\item
  Умение быстро и качественно оформлять отчёты необходимо в учебной и
  научной деятельности.
\item
  Инструменты (pandoc) позволяют автоматически преобразовывать Markdown
  в различные форматы (pdf, docx, html).
\end{itemize}
\end{frame}

\begin{frame}{Объект и предмет исследования}
\phantomsection\label{ux43eux431ux44aux435ux43aux442-ux438-ux43fux440ux435ux434ux43cux435ux442-ux438ux441ux441ux43bux435ux434ux43eux432ux430ux43dux438ux44f}
\begin{itemize}[<+->]
\tightlist
\item
  \textbf{Объект:} процесс оформления отчётов по лабораторным работам.
\item
  \textbf{Предмет:} язык разметки Markdown и инструменты его обработки
  (Pandoc).
\end{itemize}
\end{frame}

\section{Научная новизна и практическая
значимость}\label{ux43dux430ux443ux447ux43dux430ux44f-ux43dux43eux432ux438ux437ux43dux430-ux438-ux43fux440ux430ux43aux442ux438ux447ux435ux441ux43aux430ux44f-ux437ux43dux430ux447ux438ux43cux43eux441ux442ux44c}

\begin{frame}{Научная новизна}
\phantomsection\label{ux43dux430ux443ux447ux43dux430ux44f-ux43dux43eux432ux438ux437ux43dux430}
\begin{itemize}[<+->]
\tightlist
\item
  Освоение современных инструментов автоматизации создания отчётов.
\item
  Применение подхода «единый источник» для генерации документов в разных
  форматах.
\end{itemize}
\end{frame}

\begin{frame}{Практическая значимость}
\phantomsection\label{ux43fux440ux430ux43aux442ux438ux447ux435ux441ux43aux430ux44f-ux437ux43dux430ux447ux438ux43cux43eux441ux442ux44c}
\begin{itemize}[<+->]
\tightlist
\item
  Полученные навыки позволяют существенно ускорить подготовку отчётов по
  другим дисциплинам.
\item
  Возможность создавать документы, которые легко конвертировать в
  требуемый формат.
\end{itemize}
\end{frame}

\section{Цель и
задачи}\label{ux446ux435ux43bux44c-ux438-ux437ux430ux434ux430ux447ux438}

\begin{frame}{Цель}
\phantomsection\label{ux446ux435ux43bux44c}
Научиться оформлять отчёты с помощью легковесного языка разметки
Markdown.
\end{frame}

\begin{frame}{Задачи}
\phantomsection\label{ux437ux430ux434ux430ux447ux438}
\begin{itemize}[<+->]
\tightlist
\item
  Подготовить отчёт по предыдущей лабораторной работе в формате
  Markdown.
\item
  Скомпилировать отчёт в форматы pdf, docx и предоставить архив с
  исходниками.
\end{itemize}
\end{frame}

\section{Материалы и
методы}\label{ux43cux430ux442ux435ux440ux438ux430ux43bux44b-ux438-ux43cux435ux442ux43eux434ux44b}

\begin{frame}[fragile]{Материалы и методы}
\phantomsection\label{ux43cux430ux442ux435ux440ux438ux430ux43bux44b-ux438-ux43cux435ux442ux43eux434ux44b-1}
\begin{itemize}[<+->]
\tightlist
\item
  \textbf{Инструменты:}

  \begin{itemize}[<+->]
  \tightlist
  \item
    Язык разметки \textbf{Markdown}.
  \item
    Конвертер документов \textbf{Pandoc}.
  \item
    Фильтры: \texttt{pandoc-crossref} для нумерации рисунков и формул.
  \item
    Система сборки \textbf{Make} (Makefile).
  \item
    Целевые форматы: \textbf{pdf} (через LaTeX) и \textbf{docx}.
  \end{itemize}
\end{itemize}
\end{frame}

\section{Ход работы --- оформление отчёта в Markdown
(этапы)}\label{ux445ux43eux434-ux440ux430ux431ux43eux442ux44b-ux43eux444ux43eux440ux43cux43bux435ux43dux438ux435-ux43eux442ux447ux451ux442ux430-ux432-markdown-ux44dux442ux430ux43fux44b}

\begin{frame}{Этап 1: Создание шаблона отчёта}
\phantomsection\label{ux44dux442ux430ux43f-1-ux441ux43eux437ux434ux430ux43dux438ux435-ux448ux430ux431ux43bux43eux43dux430-ux43eux442ux447ux451ux442ux430}
\begin{itemize}[<+->]
\tightlist
\item
  Создание шаблона отчёта в Markdown с титульным листом, теоретическим
  введением, ходом работы и выводами.
\end{itemize}
\end{frame}

\begin{frame}{Этап 2: Вставка скриншотов}
\phantomsection\label{ux44dux442ux430ux43f-2-ux432ux441ux442ux430ux432ux43aux430-ux441ux43aux440ux438ux43dux448ux43eux442ux43eux432}
\begin{itemize}[<+->]
\tightlist
\item
  Вставка скриншотов, демонстрирующих процесс выполнения лабораторной
  работы.
\end{itemize}
\end{frame}

\begin{frame}{Этап 3: Ответы на контрольные вопросы}
\phantomsection\label{ux44dux442ux430ux43f-3-ux43eux442ux432ux435ux442ux44b-ux43dux430-ux43aux43eux43dux442ux440ux43eux43bux44cux43dux44bux435-ux432ux43eux43fux440ux43eux441ux44b}
\begin{itemize}[<+->]
\tightlist
\item
  Оформление ответов на контрольные вопросы в Markdown.
\end{itemize}
\end{frame}

\begin{frame}{Этап 4: Компиляция отчёта}
\phantomsection\label{ux44dux442ux430ux43f-4-ux43aux43eux43cux43fux438ux43bux44fux446ux438ux44f-ux43eux442ux447ux451ux442ux430}
\begin{itemize}[<+->]
\tightlist
\item
  Компиляция отчёта в pdf и docx с помощью Pandoc.
\end{itemize}
\end{frame}

\begin{frame}{Этап 5: Автоматизация сборки}
\phantomsection\label{ux44dux442ux430ux43f-5-ux430ux432ux442ux43eux43cux430ux442ux438ux437ux430ux446ux438ux44f-ux441ux431ux43eux440ux43aux438}
\begin{itemize}[<+->]
\tightlist
\item
  Автоматизация процесса с помощью Makefile.
\end{itemize}
\end{frame}

\section{Базовые элементы
Markdown}\label{ux431ux430ux437ux43eux432ux44bux435-ux44dux43bux435ux43cux435ux43dux442ux44b-markdown}

\begin{frame}[fragile]{Основные конструкции}
\phantomsection\label{ux43eux441ux43dux43eux432ux43dux44bux435-ux43aux43eux43dux441ux442ux440ux443ux43aux446ux438ux438}
\begin{itemize}[<+->]
\tightlist
\item
  \textbf{Заголовки:} \texttt{\#}, \texttt{\#\#}, \texttt{\#\#\#}
\item
  \textbf{Выделение текста:} \texttt{**жирный**}, \texttt{*курсив*},
  \texttt{***жирный\ курсив***}
\item
  \textbf{Списки:} маркированные (\texttt{-}) и нумерованные
  (\texttt{1.})
\item
  \textbf{Ссылки:} \texttt{{[}текст{]}(ссылка)}
\item
  \textbf{Вставка изображений:} \texttt{!{[}Подпись{]}(путь)}
\item
  \textbf{Формулы:} внутристрочные (\texttt{\$...\$}) и выключенные
  (\texttt{\$\$...\$\$})
\item
  \textbf{Блоки кода:} с подсветкой синтаксиса
\end{itemize}
\end{frame}

\begin{frame}{Пример использования формул}
\phantomsection\label{ux43fux440ux438ux43cux435ux440-ux438ux441ux43fux43eux43bux44cux437ux43eux432ux430ux43dux438ux44f-ux444ux43eux440ux43cux443ux43b}
\[
\sin^2 (x) + \cos^2 (x) = 1
\] \{\#eq:eq:sin2+cos2\}

Ссылка на формулу ({[}-@eq:eq:sin2+cos2{]}).
\end{frame}

\section{Скриншоты процесса
выполнения}\label{ux441ux43aux440ux438ux43dux448ux43eux442ux44b-ux43fux440ux43eux446ux435ux441ux441ux430-ux432ux44bux43fux43eux43bux43dux435ux43dux438ux44f}

\begin{frame}{Скриншоты отчёта (часть 1)}
\phantomsection\label{ux441ux43aux440ux438ux43dux448ux43eux442ux44b-ux43eux442ux447ux451ux442ux430-ux447ux430ux441ux442ux44c-1}
\includegraphics[width=0.3\linewidth,height=\textheight,keepaspectratio]{image/1.png}
\includegraphics[width=0.3\linewidth,height=\textheight,keepaspectratio]{image/2.png}
\end{frame}

\begin{frame}{Скриншоты отчёта (часть 2)}
\phantomsection\label{ux441ux43aux440ux438ux43dux448ux43eux442ux44b-ux43eux442ux447ux451ux442ux430-ux447ux430ux441ux442ux44c-2}
\includegraphics[width=0.3\linewidth,height=\textheight,keepaspectratio]{image/3.png}
\includegraphics[width=0.3\linewidth,height=\textheight,keepaspectratio]{image/4.png}
\end{frame}

\begin{frame}{Скриншоты отчёта (часть 3)}
\phantomsection\label{ux441ux43aux440ux438ux43dux448ux43eux442ux44b-ux43eux442ux447ux451ux442ux430-ux447ux430ux441ux442ux44c-3}
\includegraphics[width=0.3\linewidth,height=\textheight,keepaspectratio]{image/5.png}
\includegraphics[width=0.3\linewidth,height=\textheight,keepaspectratio]{image/6.png}
\includegraphics[width=0.3\linewidth,height=\textheight,keepaspectratio]{image/7.png}
\end{frame}

\section{Результаты}\label{ux440ux435ux437ux443ux43bux44cux442ux430ux442ux44b}

\begin{frame}[fragile]{Makefile для автоматизации сборки}
\phantomsection\label{makefile-ux434ux43bux44f-ux430ux432ux442ux43eux43cux430ux442ux438ux437ux430ux446ux438ux438-ux441ux431ux43eux440ux43aux438}
\begin{verbatim}
FILES = $(patsubst %.md, %.docx, $(wildcard *.md))
FILES += $(patsubst %.md, %.pdf, $(wildcard *.md))

LATEX_FORMAT =

FILTER = --filter pandoc-crossref

%.docx: %.md
pandoc "$<" $(FILTER) -o "$@"

%.pdf: %.md
pandoc "$<" $(LATEX_FORMAT) $(FILTER) -o "$@"

all: $(FILES)
@echo $(FILES)

clean:
rm -f $(FILES) *~
\end{verbatim}
\end{frame}

\section{Анализ и практическая
значимость}\label{ux430ux43dux430ux43bux438ux437-ux438-ux43fux440ux430ux43aux442ux438ux447ux435ux441ux43aux430ux44f-ux437ux43dux430ux447ux438ux43cux43eux441ux442ux44c}

\begin{frame}{Анализ и практическая значимость}
\begin{itemize}[<+->]
\tightlist
\item
  Создан полноценный отчёт по предыдущей лабораторной работе в Markdown.
\item
  Освоены приёмы вставки изображений, формул, листингов.
\item
  Настроена автоматическая сборка отчёта в pdf и docx.
\item
  Получены навыки, необходимые для подготовки курсовых и дипломных
  работ.
\end{itemize}
\end{frame}

\section{Выводы}\label{ux432ux44bux432ux43eux434ux44b}

\begin{frame}{Выводы}
\phantomsection\label{ux432ux44bux432ux43eux434ux44b-1}
\begin{itemize}[<+->]
\tightlist
\item
  В ходе лабораторной работы я научился оформлять отчёты с помощью языка
  разметки Markdown.
\item
  Познакомился с возможностями Pandoc по конвертации Markdown-документов
  в различные форматы.
\item
  Приобрёл опыт использования Makefile для автоматизации сборки.
\item
  Полученные компетенции будут полезны в дальнейшей учебной и научной
  деятельности.
\end{itemize}
\end{frame}




\end{document}
